\section{Protected integration and the orientation character: attack--heal}
\label{sec:agent15d-protected-integration-orientation-character}

\providecommand{\Z}{\mathbb Z}
\providecommand{\C}{\mathbb C}
\providecommand{\La}{\Lambda}
\providecommand{\Ran}{\operatorname{Ran}}
\providecommand{\RHom}{\operatorname{RHom}}
\providecommand{\sdet}{\operatorname{sdet}}
\providecommand{\CoHA}{\operatorname{CoHA}}
\providecommand{\Hall}{\mathcal H}
\providecommand{\BM}{\operatorname{BM}}
\providecommand{\Crit}{\operatorname{Crit}}
\providecommand{\MHM}{\operatorname{MHM}}
\providecommand{\Perf}{\operatorname{Perf}}
\providecommand{\Prim}{\operatorname{Prim}}
\providecommand{\sdim}{\operatorname{sdim}}
\providecommand{\Spch}{\mathrm{Sp}^{\mathrm{ch}}}
\providecommand{\hCS}{\mathrm{hCS}}
\providecommand{\BV}{\mathrm{BV}}
\providecommand{\TS}{\mathrm{TS}}
\providecommand{\Deltafive}{\Delta_5}
\providecommand{\nuDelta}{\nu_{\Delta_5}}
\providecommand{\epsor}{\epsilon_o}
\providecommand{\gDelta}{\mathfrak g_{\Delta_5}}

\subsection{Attacked claim}

The attacked claim is the strongest compact-realization upgrade:
\[
  \hbox{generic CY}_3\hbox{ orientation technology}
  + \hbox{OP Behrend integration}
  + \hbox{Borcherds denominator}
  \Longrightarrow
  \hbox{protected integration for }K3\times E
  \hbox{ with }\epsor=\nuDelta .
\]
This implication is false.

The scalar OP/DT branch is theorem-level after importing
Oberdieck--Pixton, Oberdieck--Pandharipande, and Oberdieck.  The
automorphic character
\[
  \nuDelta:\Sp_4(\Z)\to\{\pm1\}
\]
and its type-II Weyl restriction are theorem-level in the Igusa paper.
Generic orientability of CY\(_3\) moduli stacks is theorem-level in the
Joyce/Kontsevich--Soibelman/PTVV/Kinjo--Park--Safronov direction, as
used in the surrounding volumes.  None of these statements constructs:
\[
  I^{\mathrm{prot}}_{\sigma,S,o}:
  \widehat{\Hall}^{\mathrm{or}}_{K3\times E,\sigma,S}
  \longrightarrow
  \widehat{\mathbb T}^{\mathrm{prot}}_{\Gamma,S,o},
\]
a lifted type-II Weyl action on its oriented source, or the
determinant-line monodromy calculation producing
\[
  \epsor(w)=\nuDelta(w).
\]

The healed statement is conditional.

\begin{quote}
If a finite-first oriented compact critical Hall/CoHA atlas for the
relevant \(K3\times E\) sector is constructed, if its vanishing-cycle
coefficient systems admit a protected integration map compatible with
Hall multiplication, wall crossing, HN completion, and primitive
projection, and if the type-II Weyl group lifts to the oriented
completed Hall object, then the equality
\(\epsor=\nuDelta\) becomes a concrete determinant-line monodromy
problem.  It is proved only after the three simple real-root reflections
are computed geometrically and the Coxeter relations are checked.
\end{quote}

\subsection{Status ledger}

\begin{center}
\small
\begin{tabular}{p{0.34\textwidth}p{0.14\textwidth}p{0.42\textwidth}}
\hline
Claim & Status & Reason \\
\hline
\(\nuDelta\) and its type-II Weyl restriction &
theorem &
Maass character and exterior-square/reflection calculation in
\texttt{main.tex:620--760}, \texttt{main.tex:1404--1786}. \\
Reduced scalar branch
\(-4096\Deltafive^{-2}\) &
imported theorem &
OP/DT reduced quotient integration; see
\texttt{main.tex:824--1165}. \\
Scalar Behrend quotient integration &
theorem/imported scalar construction &
It lands in numbers coefficientwise.  It is not a Hall product and not
an orientation character; see \texttt{main.tex:854--888}. \\
Generic CY\(_3\) orientability of the derived moduli stack &
theorem elsewhere &
The surrounding Vol. III text uses PTVV/Joyce/KPS-type orientation data;
see
\texttt{calabi-yau-quantum-groups/chapters/examples/k3\_chiral\_bialgebra\_platonic.tex:4130--4205}. \\
Finite DWR/Ran oriented hCS--Hall comparison &
conditional theorem &
Vol. III proves sufficiency after the finite primitive package
\((\eta_{\mathrm{MC}},\lambda_{\mathrm{or}},\eta_{\mathrm{gr}},
H_{\mathrm{TS}},H_{\mathrm{fact}},Q)\) is supplied; see
\texttt{cy3\_chain\_level\_bridge.tex:3554--3695}. \\
Compact protected integration for \(K3\times E\) &
open &
The compact atlas, HN tower, proper pushforwards, reduced quotient
descent, wrapped-sector compatibility, and primitive projection are not
constructed. \\
\(\epsor=\nuDelta\) &
open &
The automorphic side exists.  The Hall-side Weyl lift and
determinant-line monodromy computation do not. \\
\hline
\end{tabular}
\end{center}

\subsection{Definition candidates}

\paragraph{Candidate A: scalar Behrend quotient.}
For Hilbert schemes or stable-pair spaces the manuscript uses
\[
  \int_{\mathrm{Hilb}^n(X,(\beta,d))/E}\nu_{\mathrm{Hilb}}\,de,
  \qquad
  \int_{P_n(X,(\beta,d))/E}\nu_P\,de .
\]
This is a coefficientwise scalar Euler characteristic.  It is the right
shadow for the OP/DT theorem.  It is not protected Hall integration:
there is no extension correspondence, no multiplication, no
orientation-line transport, no primitive projection, and no Weyl action.

\paragraph{Candidate B: motivic/critical CoHA integration.}
Let
\[
  \mathfrak M^\sigma_\gamma
\]
be the semistable compact critical/d-critical moduli stack of charge
\(\gamma\), in an extension-closed sector \(S\).  Assume finite HN
truncations \(L_{\le R}\), vanishing-cycle sheaves
\(\phi_{\mathcal W_\gamma}\), and a strong orientation \(o\) with
orientation local systems \(\mathscr L_{o,\gamma}\).  The finite-stage
oriented Hall group is
\[
  \Hall^{\mathrm{or}}_{\sigma,S,\le R}
  =
  \bigoplus_{\gamma\in L_{\le R}}
  H^{\BM}_{G_\gamma}
  \left(
    \mathfrak M^\sigma_\gamma,\,
    \phi_{\mathcal W_\gamma}\otimes\mathscr L_{o,\gamma}
  \right).
\]
The completed source is
\[
  \widehat{\Hall}^{\mathrm{or}}_{\sigma,S}
  =
  \varprojlim_R \Hall^{\mathrm{or}}_{\sigma,S,\le R}.
\]
Hall multiplication is the pull--push operation along
\[
  \mathfrak M^\sigma_\gamma\times\mathfrak M^\sigma_\eta
  \xleftarrow{\ p_{\gamma,\eta}\ }
  \mathfrak E^\sigma_{\gamma,\eta}
  \xrightarrow{\ q_{\gamma,\eta}\ }
  \mathfrak M^\sigma_{\gamma+\eta},
\]
with \(\TS\) for vanishing cycles and orientation-line transport.
The map \(q_{\gamma,\eta}\) must be proper on each finite HN quotient,
or the compactly supported pushforward is not defined.

The protected integration map should be a continuous algebra map
\[
  I^{\mathrm{prot}}_{\sigma,S,o}:
  \widehat{\Hall}^{\mathrm{or}}_{\sigma,S}
  \longrightarrow
  \widehat{\mathbb T}^{\mathrm{prot}}_{\Gamma,S,o},
\]
where the protected quantum torus has basis \(x^\gamma\) and product
\[
  x^\gamma\star_o x^\eta
  =
  \mathbb L^{-\langle\gamma,\eta\rangle/2}
  c_o(\gamma,\eta)x^{\gamma+\eta}.
\]
Here \(c_o\) is the orientation quadratic-refinement cocycle.  Its
Euler/Hodge specialization is the signed protected trace.  The
compatibility equation is
\[
  I^{\mathrm{prot}}_{\sigma,S,o}(a*b)
  =
  I^{\mathrm{prot}}_{\sigma,S,o}(a)\star_o
  I^{\mathrm{prot}}_{\sigma,S,o}(b).
\]
This is the first type-correct definition candidate.

\paragraph{Candidate C: Hodge or spin-refined protected realization.}
For a refined theory one can replace the Euler specialization by a
mixed-Hodge or perverse-filtration realization:
\[
  I^{\mathrm{prot,Hdg}}_{\sigma,S,o}:
  \widehat{\Hall}^{\mathrm{or}}_{\sigma,S}
  \longrightarrow
  \widehat{K_0(\MHM)[\Gamma]}.
\]
The Igusa paper currently needs only the unrefined protected
supertrace.  A refined target is useful only if its specialization gives
the same unrefined coefficients \(f(nm,l)\) and its orientation
transport descends to the same Weyl character.

\paragraph{Candidate D: primitive integration.}
The expression
\[
  I^{\mathrm{prot}}:\Hall\to \gDelta
\]
is not a definition of protected integration.  It confuses integration
with primitive recognition.  The Hall product and coproduct first define
primitives; protected integration gives their signed numerical shadow.
The isomorphism with \(\gDelta\) is a presentation-level theorem about
the primitive bracket, not the codomain of integration.

\subsection{Exact domain and codomain for the orientation character}

There are two different signs.

\paragraph{Hall product cocycle.}
The cocycle \(c_o(\gamma,\eta)\) is a two-variable sign or motivic
half-Tate factor attached to orientation transport for an extension
\(\gamma\to\gamma+\eta\to\eta\).  It lives in the multiplication law of
\(\widehat{\mathbb T}^{\mathrm{prot}}_{\Gamma,S,o}\).

\paragraph{Weyl monodromy character.}
The desired character is one-variable:
\[
  \epsor:
  W^{(2)}(\La^{2,1}_{II})\longrightarrow\{\pm1\}.
\]
It is not defined until there is a lifted action
\[
  \rho_o:
  W^{(2)}(\La^{2,1}_{II})
  \longrightarrow
  \operatorname{Aut}
  \bigl(\widehat{\Hall}^{\mathrm{or}}_{\sigma,S}\bigr)
\]
or, more honestly before chamber identification, an action on the
Hall wall-crossing groupoid
\[
  \rho_o(w):
  \widehat{\Hall}^{\mathrm{or}}_{\sigma,S}
  \xrightarrow{\sim}
  \widehat{\Hall}^{\mathrm{or}}_{w\sigma,wS}.
\]
To get a character on a fixed base chamber one must also supply a
wall-crossing transport
\[
  \mathsf{WC}_w:
  \widehat{\Hall}^{\mathrm{or}}_{w\sigma,wS}
  \xrightarrow{\sim}
  \widehat{\Hall}^{\mathrm{or}}_{\sigma,S}
\]
compatible with protected integration.  Then the sign is the ratio by
which \(\mathsf{WC}_w\circ\rho_o(w)\) transports the chosen orientation
branch:
\[
  (\mathsf{WC}_w\circ\rho_o(w))^*\ell_{w\gamma}
  \simeq
  \epsor(w)\,\ell_{\gamma}.
\]
The codomain is \(\{\pm1\}\) only after the ratio is central and
constant on the relevant finite HN tower.  Multiplicativity
\(\epsor(w_1w_2)=\epsor(w_1)\epsor(w_2)\) is a theorem to prove; it is
not a definition.

\subsection{Monodromy computation requirements}

The target equality is
\[
  \epsor
  =
  \nuDelta|_{W^{(2)}(\La^{2,1}_{II})}.
\]
The automorphic side is known.  Since the type-II Weyl group is
generated by the simple reflections in the real roots
\[
  \delta_1=2f_2-f_3,\qquad
  \delta_2=2f_{-2}-f_3,\qquad
  \delta_3=f_3,
\]
it is enough, after the group action and relation check exist, to prove
\[
  \epsor(s_{\delta_1})
  =
  \epsor(s_{\delta_2})
  =
  \epsor(s_{\delta_3})
  =
  -1.
\]

The geometric computation has the following irreducible obligations.

\begin{enumerate}
\item \textbf{Wall identification.}  Identify the Hall stability wall
crossed by \(s_{\delta_i}\) with the automorphic diagonal real-root
wall used by the exterior-square embedding into
\(\Sp_4(\Z)/\{\pm I\}\).  This includes the charge map
\[
  \alpha(n,l,m)=2nf_2-lf_3+2mf_{-2}.
\]
\item \textbf{Local d-critical normal form.}  Give a d-critical or
critical-chart model near the simple real-root stratum.  The model must
separate tangential modes from the single normal mode responsible for
the reflection.
\item \textbf{Orientation-line decomposition.}  Decompose
\[
  \ell_\gamma^{\otimes2}
  \simeq
  \sdet\RHom(\mathcal E_\gamma,\mathcal E_\gamma)
\]
into tangential and normal factors, with the same decomposition on the
reflected charge.
\item \textbf{Pfaffian sign.}  Prove that the reflection reverses
exactly one oriented normal square-root factor and that all tangential,
paired, and Tate-shift contributions cancel.  This is the expected
Pfaffian mechanism.  It is not computed in the current compact theory.
\item \textbf{Vanishing-cycle compatibility.}  Verify that the
monodromy of the orientation local system commutes with the vanishing
cycle functor and with \(\TS\) on the extension correspondence.
\item \textbf{Compact-support pushforward.}  Check that the wall
transport and Hall pushforwards remain proper on every finite HN
quotient and are continuous under the inverse limit.  The
\(\varprojlim^1\) obstruction named in Vol. III must vanish.
\item \textbf{Coxeter relations.}  After computing the signs on simple
reflections, verify that the transported orientation branches satisfy
the type-II Weyl relations.  Without this, the signs are local wall
data, not a character.
\item \textbf{Central lift audit.}  Compare the Hall lift with the
automorphic exterior-square lift.  A hidden metaplectic or orientation
central sign would change the equality with \(\nuDelta\).
\end{enumerate}

\subsection{Falsified shortcuts}

\begin{center}
\small
\begin{tabular}{p{0.31\textwidth}p{0.59\textwidth}}
\hline
Shortcut & Failure \\
\hline
Behrend-weighted OP integration is protected Hall integration &
False.  OP integration is scalar and coefficientwise.  It has no Hall
product, no primitive projection, and no orientation-monodromy action. \\
\([q^{1/2}r^{1/2}s^{1/2}]\Deltafive=64\) fixes orientation &
False.  The factor \(64\) is the theta-leading coefficient.  Rescaling
does not change the Maass character and constructs no orientation line. \\
\(-4096\Deltafive^{-2}\) fixes the Hall sign &
False.  The square has forgotten \(\nuDelta\).  The minus sign is the
OP scalar convention. \\
Generic CY\(_3\) orientability implies \(\epsor=\nuDelta\) &
False.  Orientability supplies square roots.  It does not supply the
Weyl action or compute monodromy. \\
The Hall-side orientation torsor trivializes on a DWR cover, hence the
relative sign is fixed &
False.  Vol. III explicitly separates Hall-side torsor triviality from
the relative hCS--Hall obstruction \(o_{\mathrm{or}}\). \\
The determinant product determines primitive parity and bracket signs &
False.  The product gives signed superdimensions.  It does not determine
even/odd dimensions separately, Hall structure constants, or
Borcherds--Kac relations. \\
Quotienting by \(E\) before Hall multiplication is harmless &
False.  Relative \(E\)-position is needed for local, wrapped, and mixed
extension correspondences.  Quotient-first loses data needed by
protected integration. \\
\hline
\end{tabular}
\end{center}

\subsection{Gate-by-gate attack and heal}

\paragraph{Gate 1: compact finite-first holomorphic \(E_3\) source.}
Attack: protected integration has no source until a compact
finite-first \(E_3\) or Hall object exists.  A BV/\(\hCS\) local model
and a scalar partition function do not give a compact charge-completed
factorization algebra on \(K3\times E\).

Heal: require a supplied
\[
  \mathcal A^{E_3}_{K3\times E}
\]
with QME, anomaly cancellation, descent, charge/HN completion, and a
comparison to the compact critical Hall atlas.  Protected integration
starts after this source exists.

\paragraph{Gate 2: Ran(\(E\)) elliptic-degree obstruction.}
Attack: ordinary support-local factorization on \(\Ran(E)\) only sees
projection-finite sectors.  Positive elliptic degree projects onto all
of \(E\), so the \(s\)-direction of the Igusa product is invisible to
finite Ran support.

Heal: protected integration must be built on a hybrid local/global
source.  Local projection-finite correspondences and wrapped global
correspondences must be formed before reduced quotienting by \(E\).  The
integration map must be continuous for the mixed completion and
compatible with the local-wrapped \(\TS\) orientation transport.

\paragraph{Gate 3: oriented compact Hall/CoHA atlas.}
Attack: one cannot define \(I^{\mathrm{prot}}\) from moduli spaces
alone.  The domain requires extension-closed stacks, critical charts,
vanishing cycles, orientations, HN finite quotients, and proper Hall
pushforwards.

Heal: construct finite quotients
\[
  \Hall^{\mathrm{or}}_{\sigma,S,\le R}
\]
with saturated Hall-lower charge sets, then pass to the inverse limit.
This matches the finite-first positive-geometry language in Vol. III.
The compact \(K3\times E\) atlas remains open.

\paragraph{Gate 4: protected integration and orientation character.}
Attack: the manuscript's \(\epsor\) comparison is meaningful only after
the integration map and the oriented Weyl action exist.  The current
paper correctly marks this as open.  No prior report or cross-volume
anchor computes the determinant-line monodromy on the three real-root
walls.

Heal: define \(I^{\mathrm{prot}}_{\sigma,S,o}\) as the continuous
vanishing-cycle Hall integration map above; define \(\epsor\) as the
orientation-branch monodromy of the lifted type-II Weyl action; prove
the three simple-reflection signs by the Pfaffian normal-mode
calculation; then check Coxeter relations.

\paragraph{Gate 5: comparison \(\epsor=\nuDelta\).}
Attack: equality cannot be chosen by normalization.  The automorphic
side is a character of \(\Sp_4(\Z)\), while the Hall side is a
geometric orientation character that is not yet constructed.

Heal: use the exterior-square embedding to put both characters on
\(W^{(2)}(\La^{2,1}_{II})\).  Since \(\nuDelta(s_{\delta_i})=-1\) for
the simple type-II reflections, the comparison reduces to the geometric
sign computation above plus the relation and central-lift audits.

\paragraph{Gate 6: primitive recognition.}
Attack: protected integration can verify signed traces, but it does not
by itself prove
\[
  H^0\Prim_{\mathrm{prot}}(\mathcal F_X)/\operatorname{Rad}
  \cong
  \gDelta .
\]
The denominator fixes superdimensions, not parity splits, simple
primitive representatives, Hall bracket constants, or the radical
quotient.

Heal: after protected integration, one must still construct the
primitive bracket, parity dimensions, Chevalley/Serre/Borcherds
relations, generation, Hopf pairing, and radical quotient.  The
orientation character enters this theorem by fixing the super sign and
reflection-character compatibility, but it does not replace the
presentation proof.

\subsection{Healed conditional theorem}

\begin{theorem}[conditional protected integration and orientation character]
\label{thm:agent15d-conditional-protected-integration}
Let \(X=K3\times E\).  Assume the following data are supplied.
\begin{enumerate}
\item A compact finite-first holomorphic \(E_3\) source
\(\mathcal A^{E_3}_X\) and a chain-level specialization to a compact
factorization object \(\mathcal F_X\) over \(\Ran(E)\), including the
hybrid wrapped-sector repair.
\item An extension-closed oriented compact critical Hall/CoHA atlas
\[
  \mathfrak M^\sigma_\gamma,\quad
  \mathfrak E^\sigma_{\gamma,\eta},\quad
  \phi_{\mathcal W_\gamma},\quad
  \ell_\gamma^{\otimes2}\simeq
  \sdet\RHom(\mathcal E_\gamma,\mathcal E_\gamma)
\]
with HN finite quotients, proper Hall pushforwards, and
Thom--Sebastiani orientation coherence.
\item A continuous protected integration map
\[
  I^{\mathrm{prot}}_{\sigma,S,o}:
  \widehat{\Hall}^{\mathrm{or}}_{X,\sigma,S}
  \to
  \widehat{\mathbb T}^{\mathrm{prot}}_{\Gamma,S,o}
\]
compatible with Hall product, wall crossing, HN completion, primitive
projection, and reduced \(E\)-quotient only after the Hall
correspondences are formed.
\item A lifted type-II Weyl action on the oriented Hall wall-crossing
groupoid, preserving the charge map
\[
  \alpha(n,l,m)=2nf_2-lf_3+2mf_{-2}
\]
and compatible with \(I^{\mathrm{prot}}_{\sigma,S,o}\).
\item For the three simple real-root reflections
\[
  s_{\delta_1},\quad s_{\delta_2},\quad s_{\delta_3},
  \qquad
  \delta_1=2f_2-f_3,\quad
  \delta_2=2f_{-2}-f_3,\quad
  \delta_3=f_3,
\]
the determinant-line normal-mode calculation gives
\[
  \epsor(s_{\delta_i})=-1.
\]
\item The orientation transports satisfy the Coxeter relations and the
Hall lift agrees with the exterior-square automorphic lift without an
extra central sign.
\end{enumerate}
Then the lifted Hall orientation monodromy defines a character
\[
  \epsor:W^{(2)}(\La^{2,1}_{II})\to\{\pm1\}
\]
and
\[
  \epsor
  =
  \nuDelta|_{W^{(2)}(\La^{2,1}_{II})}.
\]
\end{theorem}

\begin{proof}
The first three hypotheses define the oriented completed Hall algebra
and the continuous protected integration map.  The fourth hypothesis
lets the type-II Weyl group act on the oriented wall-crossing groupoid,
so the ratio of transported orientation square roots is defined.  The
fifth hypothesis computes this ratio as \(-1\) on the simple
reflections.  The sixth hypothesis makes the local signs into a
character and identifies the group element with the automorphic
reflection.  The Igusa manuscript computes
\(\nuDelta(s_{\delta_i})=-1\) on the same simple reflections.  Since
both characters agree on the generators and satisfy the same relations,
they agree on \(W^{(2)}(\La^{2,1}_{II})\).
\end{proof}

Every hypothesis is load-bearing.  The theorem is not an existence
theorem for the compact datum.

\subsection{Anchors}

\begin{itemize}
\item \texttt{main.tex:57--114}: abstract separates protected index,
virtual determinant, scalar branch, and compact realization problem.
\item \texttt{main.tex:824--960}: reduced scalar quotient integration,
protected Hall integration criterion, and open orientation monodromy
problem.
\item \texttt{main.tex:1143--1165}: constants \(64\), \(4096\), and the
OP minus sign are not orientation data.
\item \texttt{main.tex:2014--2087}: positive elliptic-degree
locality obstruction and repair alternatives.
\item \texttt{main.tex:2179--2237}: compact Igusa datum, oriented Hall
object, protected integration maps, and \(\epsor=\nuDelta\) as a
geometric condition.
\item \texttt{main.tex:2265--2361}: primitive recognition is conditional
on presentation-level Hall data.
\item \texttt{main.tex:2880--3000}: final synthesis records the three
compact consequences, including orientation monodromy.
\item \texttt{proj.bib}: Behrend 2009, Kontsevich--Soibelman CoHA, and
Davison--Meinhardt CoHA are the local bibliography anchors for
Behrend functions and critical CoHA language.
\item
\texttt{agent\_material/14\_orientations\_protected\_integration\_from\_volumes\_attack\_heal.tex}:
prior orientation report; it isolates the need for \(\rho_o\) and the
simple-reflection sign computation.
\item
\texttt{agent\_material/14\_ranE\_hybrid\_wrapping\_from\_volumes\_attack\_heal.tex}:
ordinary Ran locality misses positive elliptic degree.
\item
\texttt{agent\_material/14\_primitive\_borcherds\_recognition\_from\_volumes\_attack\_heal.tex}:
primitive recognition needs parity, relations, Hopf pairing, and
radical quotient.
\item
\texttt{calabi-yau-quantum-groups/chapters/theory/cy3\_chain\_level\_bridge.tex:522--700}:
oriented hCS--Hall comparison datum and sufficiency theorem.
\item
\texttt{calabi-yau-quantum-groups/chapters/theory/cy3\_chain\_level\_bridge.tex:1099--1211}:
relative obstruction tuple
\((o_{\mathrm{MC}},o_{\mathrm{or}},o_{\mathrm{gr}},
o_{\mathrm{TS}},o_{\mathrm{fact}})\).
\item
\texttt{calabi-yau-quantum-groups/chapters/theory/cy3\_chain\_level\_bridge.tex:3452--3695}:
Hall-side orientation torsor and finite DWR/Ran comparison; compact
inverse-limit compatibility remains.
\item
\texttt{calabi-yau-quantum-groups/chapters/theory/bps\_positive\_geometry\_closure.tex:1--180}:
finite-first oriented motivic Hall cosheaf and protected/motivic
integration template.
\item
\texttt{calabi-yau-quantum-groups/chapters/theory/bps\_positive\_geometry\_closure.tex:284--390}:
Igusa Hall--Borcherds quotient boundary and hCS--Hall obstruction.
\item
\texttt{calabi-yau-quantum-groups/chapters/examples/k3e\_bkm\_chapter.tex:13929--14145}:
\(-\Delta_5^{-2}\) character is numerical theorem-grade; compact CoHA
reading and full bialgebra parent are conditional.
\item
\texttt{chiral-bar-cobar/chapters/theory/theorem\_B\_scope\_platonic.tex:1150--1265}:
Humbert-wall monodromy and strict/weight-completed chiral bar--cobar
scope for the target \(\mathbf H_{\Delta_5}\).
\item
\texttt{chiral-bar-cobar/chapters/theory/theorem\_B\_scope\_platonic.tex:1618--1738}:
monodromy-refined bar filtration; useful as a target-side warning, not
a compact Hall orientation calculation.
\item
\texttt{chiral-bar-cobar-vol2/chapters/theory/factorization\_swiss\_cheese.tex:4286--4448}:
bi-based \(K3\times E\) architecture and conjectural CoHA-character
reading.
\item
\texttt{topological-strings/main.tex:2400--2490}: Costello BV package
and specialization problem; local BV graph technology does not prove the
compact Igusa Hall orientation.
\end{itemize}

\subsection{Open questions}

\begin{enumerate}
\item Which compact moduli problem supplies the correct
extension-closed Hall atlas before the reduced \(E\)-quotient:
ideal sheaves, stable pairs, a perverse heart, or a derived category
sector?
\item What is the precise finite-first HN height controlling positive
elliptic degree and mixed local/wrapped extension stacks?
\item Does the orientation line descend through the reduced \(E\)-quotient
after Hall correspondences are formed, and what is the residual
\(\Z/2\)-gerbe?
\item What is the local d-critical normal form at each type-II real-root
wall?
\item Does the reflected normal determinant line contribute exactly one
Pfaffian sign, or are there additional metaplectic, Tate, or chamber
transport signs?
\item Can the finite DWR/Ran primitive packages be chosen compatibly in
the \((N,r,L,m)\)-inverse system so that the Vol. III
\(\varprojlim^1\) obstruction vanishes?
\item Does the Hall-side Weyl lift coincide with the
\(\Sp_4(\Z)/\{\pm I\}\) exterior-square lift, or is there a central
orientation extension?
\item After \(\epsor=\nuDelta\), can the same orientation data be used in
the primitive recognition theorem without changing the Gritsenko--Nikulin
parity split?
\end{enumerate}

\subsection{Files changed and checks}

Only this file was written:
\[
  \texttt{agent\_material/15\_protected\_integration\_orientation\_character\_attack\_heal.tex}.
\]
Files read include \(\texttt{CLAUDE.md}\), \(\texttt{AGENTS.md}\),
\(\texttt{main.tex}\), \(\texttt{proj.bib}\), the six
\(\texttt{agent\_material/14\_*}\) reports, and the cross-volume anchors
listed above.  No branch switch, checkout, stash, reset, revert, commit,
or out-of-scope edit was performed.  A repository build was not run:
this is a standalone report and is not included by \(\texttt{main.tex}\).
